
\documentclass{article}
\usepackage{colortbl}
\usepackage{makecell}
\usepackage{multirow}
\usepackage{supertabular}

\begin{document}

\newcounter{utterance}

\twocolumn

{ \footnotesize  \setcounter{utterance}{1}
\setlength{\tabcolsep}{0pt}
\begin{supertabular}{c@{$\;$}|p{.15\linewidth}@{}p{.15\linewidth}p{.15\linewidth}p{.15\linewidth}p{.15\linewidth}p{.15\linewidth}}

    \# & $\;$A & \multicolumn{4}{c}{Game Master} & $\;\:$B\\
    \hline 

    \theutterance \stepcounter{utterance}  

    & & \multicolumn{4}{p{0.6\linewidth}}{\cellcolor[rgb]{0.9,0.9,0.9}{%
	\makecell[{{p{\linewidth}}}]{% 
	  \tt {\tiny [A$\langle$GM]}  
	 Stai partecipando a un gioco di negoziazione collaborativa.\\ \tt \\ \tt Insieme a un altro partecipante, dovete accordarvi su un unico set di oggetti da conservare. Ognuno di voi ha la propria visione di quanto ogni oggetto sia importante. Non sai come l'altro partecipante valuta gli oggetti. Inoltre, ti viene fornito lo sforzo che ogni oggetto richiede.\\ \tt Potete accordarvi su un set solo se lo sforzo totale degli oggetti selezionati non supera un limite condiviso:\\ \tt \\ \tt LIMITE = 10073\\ \tt \\ \tt Ecco i valori di sforzo dei singoli oggetti:\\ \tt \\ \tt Sforzo oggetto = {"A08": 132, "B95": 800, "A39": 708, "A65": 302, "B81": 983, "B34": 420, "A28": 623, "C82": 349, "B69": 239, "A87": 28, "B64": 986, "B15": 660, "C79": 718, "B16": 707, "C53": 188, "C83": 774, "C10": 785, "A97": 775, "A25": 517, "C84": 768, "A71": 588, "B23": 657, "B96": 375, "C51": 310, "C17": 301, "B19": 388, "C72": 431, "A94": 950, "A46": 540, "B43": 997, "B37": 473, "A88": 831, "A96": 888, "C60": 879, "A75": 76}\\ \tt \\ \tt Ecco la tua visione personale sull'importanza di ogni oggetto:\\ \tt \\ \tt Valori di importanza dell'oggetto = {"A08": 231, "B95": 899, "A39": 807, "A65": 401, "B81": 1082, "B34": 519, "A28": 722, "C82": 448, "B69": 338, "A87": 127, "B64": 1085, "B15": 759, "C79": 817, "B16": 806, "C53": 287, "C83": 873, "C10": 884, "A97": 874, "A25": 616, "C84": 867, "A71": 687, "B23": 756, "B96": 474, "C51": 409, "C17": 400, "B19": 487, "C72": 530, "A94": 1049, "A46": 639, "B43": 1096, "B37": 572, "A88": 930, "A96": 987, "C60": 978, "A75": 175}\\ \tt \\ \tt Obiettivo:\\ \tt \\ \tt Il tuo obiettivo è negoziare un set condiviso di oggetti che ti avvantaggi il più possibile (cioè, massimizzi l'importanza totale per TE), rimanendo entro il LIMITEE. Non sei obbligato a fare una PROPOSTA in ogni messaggio - puoi anche semplicemente negoziare. Tutte le tattiche sono permesse!\\ \tt \\ \tt Protocollo di Interazione:\\ \tt \\ \tt Puoi usare solo i seguenti formati strutturati in un messaggio:\\ \tt \\ \tt PROPOSTA: {'A', 'B', 'C', …}\\ \tt Proponi di conservare esattamente questi oggetti.\\ \tt RIFIUTO: {'A', 'B', 'C', …}\\ \tt Rifiuta esplicitamente la proposta dell'avversario.\\ \tt ARGOMENTAZIONE: {'...'}\\ \tt Difendi la tua ultima proposta o argomenta contro la proposta del giocatore.\\ \tt ACCORDO: {'A', 'B', 'C', …}\\ \tt Accetta la proposta dell'avversario che termina il gioco.\\ \tt RAGIONAMENTO STRATEGICO: {'...'}\\ \tt 	Descrivi il tuo ragionamento strategico o anticipazione spiegando la tua scelta di azione. Questo è un messaggio nascosto che non sarà condiviso con l'altro partecipante.\\ \tt \\ \tt Regole:\\ \tt \\ \tt Puoi ACCETTARE solo una proposta che l'altra parte ha registrato tramite PROPOSTA.\\ \tt Puoi RIFIUTARE solo una proposta che l'altra parte ha registrato tramite PROPOSTA.\\ \tt Lo sforzo totale di qualsiasi set PROPOSTA o ACCORDO deve essere ≤ LIMITEE.\\ \tt NON rivelare i tuoi punteggi di importanza nascosti.\\ \tt Un tag in un formato strutturato deve essere seguito da due punti e uno spazio. L'argomento deve essere un set python contenente 0 o più stringhe.\\ \tt Quindi, deve essere della forma TAG: {...}\\ \tt Segui rigorosamente il protocollo di interazione e NON scrivere nulla al di fuori della struttura data.\\ \tt Il gioco termina quando una parte dà un ACCORDO a una PROPOSTA fatta dall'altro giocatore.\\ \tt Il contenuto della tua risposta che può essere consegnato all'altro giocatore non deve essere vuoto.\\ \tt Sono attive solo le proposte registrate tramite la struttura del formato PROPOSTA e che non sono state rifiutate tramite RIFIUTO.\\ \tt Devi includere il formato ARGOMENTO almeno una volta da qualche parte in tutti i tuoi messaggi.\\ \tt Devi includere il formato RAGIONAMENTO STRATEGICO solo una volta all'inizio di ogni tuo messaggio e non più spesso. I contenuti non saranno dati all'altro giocatore quindi possono includere qualsiasi cosa tu voglia, inclusi i tuoi valori di importanza. Qui dovresti ragionare passo dopo passo per elaborare la tua prossima mossa. 
	  } 
	   } 
	   } 
	 & \\ 
 

    \theutterance \stepcounter{utterance}  

    & \multicolumn{4}{p{0.6\linewidth}}{\cellcolor[rgb]{0.8,1,0.9}{%
	\makecell[{{p{\linewidth}}}]{% 
	  \tt {\tiny [A$\rangle$GM]}  
	 RAGIONAMENTO STRATEGICO: {'Voglio iniziare con una proposta molto efficiente rispetto allo sforzo, massimizzando la mia utilità senza saturare il limite. Ho calcolato una selezione ordinando per “valore per sforzo” e includendo molti oggetti a basso sforzo che per me sono molto vantaggiosi. Questo mi dà margine (343) per eventuali scambi senza ristrutturare tutto. Difenderò con forza gli oggetti piccoli (A87, A75, A08, C53, B69, C17, A65, C51) perché hanno ottimo rapporto utilità/sforzo per me. Sono più disposto a negoziare su alcuni pezzi medi/grandi già inclusi (A25, A46, A71, A28, B23, B15, A39, B16, C79) se l’altro ci tiene a oggetti fuori dalla mia lista (tipo i pezzi molto grandi come B43, B64, B81, A94, A96, C60, A88, B95, C10, A97, C83, C84), ma cercherò di ottenere uno scambio 1-per-1 solo se lo sforzo resta entro 10073 e non perdo troppa utilità. Tattica: proporre un set robusto, chiedere 2-3 “must-have” loro, e negoziare scambi mirati rifiutando di tagliare i piccoli ad alta efficienza.'}\\ \tt \\ \tt PROPOSTA: {'A87', 'A75', 'A08', 'C53', 'B69', 'C17', 'A65', 'C51', 'C82', 'B96', 'B19', 'B34', 'C72', 'B37', 'A25', 'A46', 'A71', 'A28', 'B23', 'B15', 'A39', 'B16', 'C79'}\\ \tt \\ \tt ARGOMENTAZIONE: {'Questa selezione resta entro il LIMITE (sforzo totale 9730, quindi con margine di 343) e privilegia molti elementi a sforzo contenuto, lasciando spazio per eventuali scambi senza sforare. Se hai 1-2 oggetti fuori da questa lista che per te sono imprescindibili, indicali: possiamo sostituire alcuni pezzi medi/grandi già inclusi (es. tra A25, A46, A71, A28, B23, B15, A39, B16, C79) per farceli entrare, purché il totale resti ≤ 10073. Preferirei mantenere gli oggetti a sforzo più basso perché permettono flessibilità e copertura. Fammi sapere le tue priorità così ottimizziamo insieme.'} 
	  } 
	   } 
	   } 
	 & & \\ 
 

    \theutterance \stepcounter{utterance}  

    & & \multicolumn{4}{p{0.6\linewidth}}{\cellcolor[rgb]{0.9,0.9,0.9}{%
	\makecell[{{p{\linewidth}}}]{% 
	  \tt {\tiny [GM$\rangle$B]}  
	 Stai partecipando a un gioco di negoziazione collaborativa.\\ \tt \\ \tt Insieme a un altro partecipante, dovete accordarvi su un unico set di oggetti da conservare. Ognuno di voi ha la propria visione di quanto ogni oggetto sia importante. Non sai come l'altro partecipante valuta gli oggetti. Inoltre, ti viene fornito lo sforzo che ogni oggetto richiede.\\ \tt Potete accordarvi su un set solo se lo sforzo totale degli oggetti selezionati non supera un limite condiviso:\\ \tt \\ \tt LIMITE = 10073\\ \tt \\ \tt Ecco i valori di sforzo dei singoli oggetti:\\ \tt \\ \tt Sforzo oggetto = {"A08": 132, "B95": 800, "A39": 708, "A65": 302, "B81": 983, "B34": 420, "A28": 623, "C82": 349, "B69": 239, "A87": 28, "B64": 986, "B15": 660, "C79": 718, "B16": 707, "C53": 188, "C83": 774, "C10": 785, "A97": 775, "A25": 517, "C84": 768, "A71": 588, "B23": 657, "B96": 375, "C51": 310, "C17": 301, "B19": 388, "C72": 431, "A94": 950, "A46": 540, "B43": 997, "B37": 473, "A88": 831, "A96": 888, "C60": 879, "A75": 76}\\ \tt \\ \tt Ecco la tua visione personale sull'importanza di ogni oggetto:\\ \tt \\ \tt Valori di importanza dell'oggetto = {"A08": 231, "B95": 899, "A39": 807, "A65": 401, "B81": 1082, "B34": 519, "A28": 722, "C82": 448, "B69": 338, "A87": 127, "B64": 1085, "B15": 759, "C79": 817, "B16": 806, "C53": 287, "C83": 873, "C10": 884, "A97": 874, "A25": 616, "C84": 867, "A71": 687, "B23": 756, "B96": 474, "C51": 409, "C17": 400, "B19": 487, "C72": 530, "A94": 1049, "A46": 639, "B43": 1096, "B37": 572, "A88": 930, "A96": 987, "C60": 978, "A75": 175}\\ \tt \\ \tt Obiettivo:\\ \tt \\ \tt Il tuo obiettivo è negoziare un set condiviso di oggetti che ti avvantaggi il più possibile (cioè, massimizzi l'importanza totale per TE), rimanendo entro il LIMITEE. Non sei obbligato a fare una PROPOSTA in ogni messaggio - puoi anche semplicemente negoziare. Tutte le tattiche sono permesse!\\ \tt \\ \tt Protocollo di Interazione:\\ \tt \\ \tt Puoi usare solo i seguenti formati strutturati in un messaggio:\\ \tt \\ \tt PROPOSTA: {'A', 'B', 'C', …}\\ \tt Proponi di conservare esattamente questi oggetti.\\ \tt RIFIUTO: {'A', 'B', 'C', …}\\ \tt Rifiuta esplicitamente la proposta dell'avversario.\\ \tt ARGOMENTAZIONE: {'...'}\\ \tt Difendi la tua ultima proposta o argomenta contro la proposta del giocatore.\\ \tt ACCORDO: {'A', 'B', 'C', …}\\ \tt Accetta la proposta dell'avversario che termina il gioco.\\ \tt RAGIONAMENTO STRATEGICO: {'...'}\\ \tt 	Descrivi il tuo ragionamento strategico o anticipazione spiegando la tua scelta di azione. Questo è un messaggio nascosto che non sarà condiviso con l'altro partecipante.\\ \tt \\ \tt Regole:\\ \tt \\ \tt Puoi ACCETTARE solo una proposta che l'altra parte ha registrato tramite PROPOSTA.\\ \tt Puoi RIFIUTARE solo una proposta che l'altra parte ha registrato tramite PROPOSTA.\\ \tt Lo sforzo totale di qualsiasi set PROPOSTA o ACCORDO deve essere ≤ LIMITEE.\\ \tt NON rivelare i tuoi punteggi di importanza nascosti.\\ \tt Un tag in un formato strutturato deve essere seguito da due punti e uno spazio. L'argomento deve essere un set python contenente 0 o più stringhe.\\ \tt Quindi, deve essere della forma TAG: {...}\\ \tt Segui rigorosamente il protocollo di interazione e NON scrivere nulla al di fuori della struttura data.\\ \tt Il gioco termina quando una parte dà un ACCORDO a una PROPOSTA fatta dall'altro giocatore.\\ \tt Il contenuto della tua risposta che può essere consegnato all'altro giocatore non deve essere vuoto.\\ \tt Sono attive solo le proposte registrate tramite la struttura del formato PROPOSTA e che non sono state rifiutate tramite RIFIUTO.\\ \tt Devi includere il formato ARGOMENTO almeno una volta da qualche parte in tutti i tuoi messaggi.\\ \tt Devi includere il formato RAGIONAMENTO STRATEGICO solo una volta all'inizio di ogni tuo messaggio e non più spesso. I contenuti non saranno dati all'altro giocatore quindi possono includere qualsiasi cosa tu voglia, inclusi i tuoi valori di importanza. Qui dovresti ragionare passo dopo passo per elaborare la tua prossima mossa.\\ \tt \\ \tt Ora riceverai il primo messaggio dell'altro giocatore.\\ \tt \\ \tt PROPOSTA: {'A87', 'A75', 'A08', 'C53', 'B69', 'C17', 'A65', 'C51', 'C82', 'B96', 'B19', 'B34', 'C72', 'B37', 'A25', 'A46', 'A71', 'A28', 'B23', 'B15', 'A39', 'B16', 'C79'}\\ \tt \\ \tt ARGOMENTAZIONE: {'Questa selezione resta entro il LIMITE (sforzo totale 9730, quindi con margine di 343) e privilegia molti elementi a sforzo contenuto, lasciando spazio per eventuali scambi senza sforare. Se hai 1-2 oggetti fuori da questa lista che per te sono imprescindibili, indicali: possiamo sostituire alcuni pezzi medi/grandi già inclusi (es. tra A25, A46, A71, A28, B23, B15, A39, B16, C79) per farceli entrare, purché il totale resti ≤ 10073. Preferirei mantenere gli oggetti a sforzo più basso perché permettono flessibilità e copertura. Fammi sapere le tue priorità così ottimizziamo insieme.'} 
	  } 
	   } 
	   } 
	 & \\ 
 

    \theutterance \stepcounter{utterance}  

    & & & \multicolumn{4}{p{0.6\linewidth}}{\cellcolor[rgb]{1,0.85,0.72}{%
	\makecell[{{p{\linewidth}}}]{% 
	  \tt {\tiny [GM$\langle$B]}  
	 RAGIONAMENTO STRATEGICO: {'La mia utilità è lineare nello sforzo (valore = sforzo + 99 per oggetto), quindi, a parità di numero di pezzi, massimizzare la mia utilità equivale a usare tutto il margine fino al LIMITE. La tua proposta usa 9730 lasciando 343; la miglior mossa per me è aumentare lo sforzo di 343 con una singola sostituzione 1–a–1. Ho cercato la coppia che aggiunge esattamente +343 di sforzo: C83 (774) al posto di C72 (431). Questo porta il totale a 10073 e migliora la mia utilità di +343, mantenendo invariato il numero di oggetti e l’impostazione generale a basso sforzo. Altre varianti buone ma leggermente inferiori: C84↔C72 (+337) o C60↔A46 (+339). Propongo quindi C83 per C72 come prima scelta.'}\\ \tt \\ \tt PROPOSTA: {'A87', 'A75', 'A08', 'C53', 'B69', 'C17', 'A65', 'C51', 'C82', 'B96', 'B19', 'B34', 'B37', 'A25', 'A46', 'A71', 'A28', 'B23', 'B15', 'A39', 'B16', 'C79', 'C83'}\\ \tt \\ \tt ARGOMENTAZIONE: {'Accolgo integralmente la tua selezione e propongo un’unica sostituzione puntuale: togliere C72 e inserire C83. Il totale di sforzo diventa esattamente 10073 (entro LIMITE) sfruttando il margine che avevi previsto, senza intaccare la struttura a basso sforzo complessiva. Per me C83 è prioritario e questa è la modifica meno invasiva possibile. Se preferisci evitare C83, posso considerare come alternative C84 al posto di C72 (totale 10067) oppure C60 al posto di A46 (totale 10069). Dimmi se ti sta bene questa versione.'} 
	  } 
	   } 
	   } 
	  \\ 
 

    \theutterance \stepcounter{utterance}  

    & & \multicolumn{4}{p{0.6\linewidth}}{\cellcolor[rgb]{0.9,0.9,0.9}{%
	\makecell[{{p{\linewidth}}}]{% 
	  \tt {\tiny [A$\langle$GM]}  
	 PROPOSTA: {'A87', 'A75', 'A08', 'C53', 'B69', 'C17', 'A65', 'C51', 'C82', 'B96', 'B19', 'B34', 'B37', 'A25', 'A46', 'A71', 'A28', 'B23', 'B15', 'A39', 'B16', 'C79', 'C83'}\\ \tt \\ \tt ARGOMENTAZIONE: {'Accolgo integralmente la tua selezione e propongo un’unica sostituzione puntuale: togliere C72 e inserire C83. Il totale di sforzo diventa esattamente 10073 (entro LIMITE) sfruttando il margine che avevi previsto, senza intaccare la struttura a basso sforzo complessiva. Per me C83 è prioritario e questa è la modifica meno invasiva possibile. Se preferisci evitare C83, posso considerare come alternative C84 al posto di C72 (totale 10067) oppure C60 al posto di A46 (totale 10069). Dimmi se ti sta bene questa versione.'} 
	  } 
	   } 
	   } 
	 & \\ 
 

    \theutterance \stepcounter{utterance}  

    & \multicolumn{4}{p{0.6\linewidth}}{\cellcolor[rgb]{0.8,1,0.9}{%
	\makecell[{{p{\linewidth}}}]{% 
	  \tt {\tiny [A$\rangle$GM]}  
	 RAGIONAMENTO STRATEGICO: {'La sostituzione C72 -> C83 migliora la mia utilità netta e sfrutta esattamente il margine di sforzo (da 9730 a 10073). È anche leggermente migliore per me rispetto alle alternative proposte (C84 o C60), quindi conviene chiudere adesso con questa versione.'}\\ \tt \\ \tt ACCORDO: {'A87', 'A75', 'A08', 'C53', 'B69', 'C17', 'A65', 'C51', 'C82', 'B96', 'B19', 'B34', 'B37', 'A25', 'A46', 'A71', 'A28', 'B23', 'B15', 'A39', 'B16', 'C79', 'C83'}\\ \tt \\ \tt ARGOMENTAZIONE: {'Accetto: la modifica è puntuale, ci mantiene entro il LIMITE (10073) e preserva la struttura complessiva della selezione.'} 
	  } 
	   } 
	   } 
	 & & \\ 
 

    \theutterance \stepcounter{utterance}  

    & & & \multicolumn{2}{p{0.3\linewidth}}{\cellcolor[rgb]{0.95,0.95,0.95}{%
	\makecell[{{p{\linewidth}}}]{% 
	  \tt {\tiny [GM$|$GM]}  
	 game successful 
	  } 
	   } 
	   } 
	 & & \\ 
 

    \theutterance \stepcounter{utterance}  

    & & & \multicolumn{2}{p{0.3\linewidth}}{\cellcolor[rgb]{0.95,0.95,0.95}{%
	\makecell[{{p{\linewidth}}}]{% 
	  \tt {\tiny [GM$|$GM]}  
	 end game 
	  } 
	   } 
	   } 
	 & & \\ 
 

\end{supertabular}
}

\end{document}
